\section{Chỉ số như một dụng cụ đo}
\label{sec:ch08-chi-so-nhu-mot-dung-cu}

\index{chỉ số đánh giá}%
\index{đánh giá theo chức năng}%
Ba đại lượng mà chương trước để lại đều có cùng một dạng: chúng chấm điểm một
lời giải thích mà không hỏi ai cả. Doshi-Velez và Kim đã xếp sẵn dạng ấy vào một
trong ba cách đánh giá của họ. Mục~\ref{sec:ch01-danh-gia-giai-thich} đã đọc cả
ba; cách thứ ba, đánh giá theo chức năng, được các tác giả mô tả là không có
người tham gia và làm việc trên nhiệm vụ thay thế~\autocite{p14rigorous}. Mọi
chỉ số trong chương này thuộc cách ấy.

Điều đó không làm chúng kém giá trị. Một phép đo tự động chạy được hàng nghìn
lần, trên hàng nghìn đầu vào, với chi phí mà một nghiên cứu có người tham gia
không bao giờ chịu nổi, nên nếu không có nhóm chỉ số này thì phần lớn công trình
so sánh các phương pháp attribution đã không thể tồn tại. Cái giá nằm ở chỗ
khác, và Doshi-Velez và Kim nêu nó ngay trong mục định nghĩa cách thứ ba. Họ viết
rằng đánh giá theo chức năng thích hợp nhất khi lớp mô hình hoặc bộ điều chuẩn
đang xét đã được kiểm chứng từ trước, chẳng hạn bằng các thí nghiệm có người
tham gia~\autocite{p14rigorous}. Điều kiện ấy đặt cách thứ ba vào thế phụ thuộc:
nó rẻ được là vì đã có ai đó trả giá đắt trước. Còn chọn \tn{đại lượng thay thế}{proxy} nào
thì chính các tác giả để ngỏ, và họ xếp câu hỏi ấy vào danh sách vấn đề mở ở
cuối bài~\autocite{p14rigorous}.

Ta gọi tên điều kiện đó bằng một phép so sánh sẽ theo suốt chương. Một \tn{chỉ số trung thực}{faithfulness metric}
là một dụng cụ đo, và một dụng cụ đo có hai phẩm chất tách rời nhau.
Thứ nhất là nó cho ra số ổn định, lặp lại được, rẻ. Thứ hai là số ấy tương ứng
với đại lượng ghi trên mặt đồng hồ. Một nhiệt kế hỏng vẫn có thể rất ổn định.
Với các chỉ số trong chương này, phẩm chất thứ nhất được báo cáo thường xuyên,
còn phẩm chất thứ hai thì hầu như không, và bài báo neo của cuốn sách nói thẳng
điều đó: phần lớn công trình đề xuất chỉ số chỉ báo cáo điểm số tuyệt đối hoặc
so sánh với các chỉ số trước, chứ không đối chiếu với đáp án
đúng~\autocite{p21metrics}.

\begin{definition}[Chỉ số trung thực]
\label{def:ch08-chi-so-trung-thuc}
Cho mô hình $f$, đầu vào $x$ và một lời giải thích gán điểm cho các
\tn{đặc trưng}{feature} của $x$. Một chỉ số trung thực là một thủ tục nhận ba thứ ấy và trả về một số
thực, kèm quy ước rằng số càng cao thì lời giải thích càng phản ánh đúng sự phụ
thuộc của $f(x)$ vào các đặc trưng.
\end{definition}

Định nghĩa~\ref{def:ch08-chi-so-trung-thuc} cố tình để hở đúng chỗ đang bàn.
Nó mô tả thủ tục và quy ước đọc kết quả, nhưng không đòi hỏi bằng chứng nào cho
quy ước ấy. Từ mục~\ref{sec:ch08-ho-xoa-dan} trở đi, mỗi họ chỉ số được kiểm lại
xem chỗ hở đó được lấp bằng gì.

Còn một khả năng nữa cần loại trừ trước, vì nó có vẻ là lối ra rẻ nhất. Một
khảo sát về XAI có thể đã chốt sẵn cách đo cho cả lĩnh vực, và khi ấy chương này
chỉ việc chép lại. Khảo sát của Long và cộng sự, mà
mục~\ref{sec:ch07-ban-do} đã đọc, không làm việc đó: nó điểm lại công cụ theo
phạm vi và theo thời điểm can thiệp, và không định nghĩa một chỉ số trung thực
nào~\autocite{p20trends}. Không có bảng nào để chép, nên các họ chỉ số phải được
dựng lại từ chính các bài báo đã đề xuất chúng.
